\documentclass[12pt, a4paper]{article}
 
% ── Packages ──────────────────────────────────────────────────────────────────
\usepackage[margin=2.5cm]{geometry}
\usepackage{amsmath, amssymb}
\usepackage{graphicx}
\usepackage{booktabs}
\usepackage{hyperref}
\usepackage{cite}
\usepackage{physics}
\usepackage{xcolor}
\usepackage{microtype}
\usepackage{setspace}
\usepackage{caption}
 
\onehalfspacing
 
\hypersetup{
    colorlinks = true,
    linkcolor  = blue!70!black,
    citecolor  = blue!70!black,
    urlcolor   = blue!70!black
}
 
% ── Title ─────────────────────────────────────────────────────────────────────
\title{%
    \textbf{Measuring the Solar Mixing Angle $\theta_{12}$}\\[4pt]
    \large A Chi-Squared Analysis of SNO Phase~I Data\\
    in the Adiabatic MSW Regime
}
 
\author{Ashok Kumar Ramu\\[2pt]
    \small Department of Physics and Astronomy, University of Waterloo}
 
\date{\today}
 
% ══════════════════════════════════════════════════════════════════════════════
\begin{document}
\maketitle
 
% ── Abstract ──────────────────────────────────────────────────────────────────
\begin{abstract}
We reproduce the measurement of the solar neutrino mixing angle $\theta_{12}$
using the charged-current (CC) and neutral-current (NC) flux measurements from
the Sudbury Neutrino Observatory (SNO) Phase~I dataset (2002).
Working in the adiabatic Mikheyev--Smirnov--Wolfenstein (MSW) regime,
appropriate for high-energy ${}^8\text{B}$ solar neutrinos,
the electron-neutrino survival probability reduces to
$P_{ee} = \sin^2\theta_{12}$.
We perform a chi-squared minimisation of the CC/NC flux ratio and extract
\[
    \theta_{12} = 36.0^{\circ} \pm 2.3^{\circ},
\]
consistent with the SNO-only result published in 2002.
We demonstrate that the vacuum-averaged formula
$P_{ee} = 1 - \tfrac{1}{2}\sin^2 2\theta_{12} \approx 0.548$
is incompatible with the observed suppression of $\sim\!34\%$,
confirming that matter effects are essential to the correct interpretation
of solar neutrino data.
The slight offset from the current PDG combined value of
$\theta_{12} = 33.4^{\circ}$ is explained qualitatively through
the Earth's day-night regeneration effect and the subsequent KamLAND
constraint on $\Delta m^2_{21}$.
\end{abstract}

\newpage 
\tableofcontents
\newpage
 
% ══════════════════════════════════════════════════════════════════════════════
\section{Introduction}
\label{sec:intro}
 
The solar neutrino problem --- the deficit of electron neutrinos arriving at
Earth relative to Standard Solar Model (SSM) predictions --- stood as one of
the most compelling anomalies in particle physics for over three decades
\cite{Homestake1968, Bahcall1989}.
Early radiochemical experiments such as Homestake, GALLEX, and SAGE observed
roughly one-third to one-half the predicted flux, but could not determine
whether the missing neutrinos had oscillated into other flavours or whether
the SSM itself was simply wrong.
 
The Sudbury Neutrino Observatory (SNO), located 2~km underground in the
INCO Creighton mine in Sudbury, Ontario (now SNOLAB), resolved this ambiguity
decisively \cite{SNO2002}.
By using heavy water ($\text{D}_2\text{O}$) as both target and Cherenkov
medium, SNO could simultaneously measure:
\begin{itemize}
    \item the \textbf{charged-current} (CC) reaction
          $\nu_e + d \to p + p + e^-$,
          sensitive only to electron neutrinos;
    \item the \textbf{neutral-current} (NC) reaction
          $\nu_x + d \to p + n + \nu_x$,
          equally sensitive to all active flavours.
\end{itemize}
The ratio $\phi_\text{CC}/\phi_\text{NC}$ directly measures the
electron-neutrino survival probability $P_{ee}$, independent of solar model
uncertainties.
When SNO reported $\phi_\text{NC} \approx \phi_\text{SSM}$ but
$\phi_\text{CC} \approx 0.35\,\phi_\text{SSM}$, it provided the first
direct, model-independent proof that solar neutrinos change flavour
in transit \cite{SNO2002}.
 
In this paper we reproduce the SNO Phase~I extraction of $\theta_{12}$,
emphasising the critical role of the MSW matter effect.
Section~\ref{sec:theory} derives the relevant survival probability from
first principles.
Section~\ref{sec:data} describes the experimental inputs and chi-squared
method.
Section~\ref{sec:results} presents our fit result and walks through each
panel of the analysis figure.
Section~\ref{sec:discussion} discusses the limitations of the SNO-only
measurement and the refinements brought by KamLAND and the Earth's
regeneration effect.
 
% ══════════════════════════════════════════════════════════════════════════════
\section{Theoretical Framework}
\label{sec:theory}
 
\subsection{Two-Flavour Neutrino Oscillations in Vacuum}
 
In vacuum, the flavour eigenstates $\ket{\nu_e}$ and $\ket{\nu_\mu}$ are
related to the mass eigenstates $\ket{\nu_1}$ and $\ket{\nu_2}$ by the
PMNS mixing matrix, which in the two-flavour approximation reduces to a
single rotation by angle $\theta_{12}$:
\begin{equation}
    \begin{pmatrix} \nu_e \\ \nu_\mu \end{pmatrix}
    =
    \begin{pmatrix}
        \cos\theta_{12} & \sin\theta_{12} \\
       -\sin\theta_{12} & \cos\theta_{12}
    \end{pmatrix}
    \begin{pmatrix} \nu_1 \\ \nu_2 \end{pmatrix}.
    \label{eq:mixing}
\end{equation}
For a neutrino propagating a distance $L$ in vacuum, the survival probability
oscillates as:
\begin{equation}
    P_{ee}^{\text{vac}}(L,E)
    = 1 - \sin^2 2\theta_{12}\,
      \sin^2\!\left(\frac{\Delta m^2_{21}\, L}{4E}\right),
    \label{eq:pvac}
\end{equation}
where $\Delta m^2_{21} = m_2^2 - m_1^2$.
Averaging over many oscillation lengths yields:
\begin{equation}
    \langle P_{ee}^{\text{vac}} \rangle
    = 1 - \tfrac{1}{2}\sin^2 2\theta_{12}.
    \label{eq:pvac_avg}
\end{equation}
 
\subsection{The MSW Effect}
\label{sec:msw}
 
Equation~\eqref{eq:pvac_avg} assumes the neutrino travels in vacuum its
entire journey.
This is incorrect for ${}^8\text{B}$ solar neutrinos, produced in the
Sun's core where the electron number density $N_e \sim 10^{26}\ \text{cm}^{-3}$.
 
In matter, the effective Hamiltonian acquires an additional potential
$V = \sqrt{2}\,G_F N_e$ for $\nu_e$, shifting the mixing angle to an
effective in-matter value $\theta_m$:
\begin{equation}
    \tan 2\theta_m
    = \frac{\sin 2\theta_{12}}
           {\cos 2\theta_{12}
            - \dfrac{2\sqrt{2}\,G_F N_e E}{\Delta m^2_{21}}}.
    \label{eq:msw_angle}
\end{equation}
The \textbf{MSW resonance} occurs when the denominator vanishes:
\begin{equation}
    N_e^{\text{res}}
    = \frac{\Delta m^2_{21}\cos 2\theta_{12}}{2\sqrt{2}\,G_F E}.
    \label{eq:resonance}
\end{equation}
For ${}^8\text{B}$ neutrinos with $E \sim 5$--$10$~MeV, this resonance
lies well within the solar interior.
 
\subsection{Adiabatic Propagation and the Survival Probability}
 
In the \textbf{adiabatic limit}, the propagation proceeds as follows:
\begin{enumerate}
    \item At production in the high-density core, $\theta_m \approx 90^\circ$,
          so the electron neutrino is born almost entirely in the heavier
          matter eigenstate $\ket{\nu_2^m}$.
    \item As the neutrino travels outward through decreasing density,
          adiabaticity ensures it remains in the same matter eigenstate.
    \item At the solar surface, $N_e \to 0$ and $\theta_m \to \theta_{12}$,
          so the neutrino exits as a pure vacuum mass eigenstate $\ket{\nu_2}$.
    \item The probability of detecting this state as $\nu_e$ is:
          \begin{equation}
              \boxed{P_{ee}^{\text{MSW}}
              = \left|\braket{\nu_e}{\nu_2}\right|^2
              = \sin^2\theta_{12}.}
              \label{eq:pmsw}
          \end{equation}
\end{enumerate}
For $\theta_{12} = 33.4^\circ$, the two formulas give strikingly
different predictions:
\begin{align}
    P_{ee}^{\text{MSW}}
        &= \sin^2(33.4^\circ) \approx 0.303, \label{eq:pmsw_num}\\
    \langle P_{ee}^{\text{vac}}\rangle
        &= 1 - \tfrac{1}{2}\sin^2(66.8^\circ) \approx 0.578.
        \label{eq:pvac_num}
\end{align}
SNO observed $P_{ee} \approx 0.346$ --- consistent with MSW and
incompatible with vacuum at the $\sim\!6\sigma$ level.
 
% ══════════════════════════════════════════════════════════════════════════════
\section{Data and Method}
\label{sec:data}
 
\subsection{SNO Phase I Flux Measurements}
 
We use the flux measurements from Reference~\cite{SNO2002}, in units of
$10^6\ \text{cm}^{-2}\,\text{s}^{-1}$:
\begin{align}
    \phi_\text{CC} &= 1.76 \pm 0.11, \label{eq:phicc}\\
    \phi_\text{NC} &= 5.09 \pm 0.45. \label{eq:phinc}
\end{align}
 
\subsection{Deriving the Experimental Survival Probability}
 
The observed electron-neutrino survival probability is:
\begin{equation}
    P_{ee}^{\text{obs}}
    = \frac{\phi_\text{CC}}{\phi_\text{NC}}
    = \frac{1.76}{5.09}
    = 0.3458,
    \label{eq:peeobs}
\end{equation}
with uncertainty propagated via:
\begin{equation}
    \frac{\sigma_{P_{ee}}}{P_{ee}}
    = \sqrt{\left(\frac{\sigma_\text{CC}}{\phi_\text{CC}}\right)^2
           +\left(\frac{\sigma_\text{NC}}{\phi_\text{NC}}\right)^2},
    \label{eq:errprop}
\end{equation}
giving:
\begin{equation}
    P_{ee}^{\text{obs}} = 0.3458 \pm 0.0374.
    \label{eq:peeobs_final}
\end{equation}
 
\subsection{Chi-Squared Fit}
 
The chi-squared function for a single-point measurement is:
\begin{equation}
    \chi^2(\theta_{12})
    = \left(
        \frac{\sin^2\theta_{12} - 0.3458}{0.0374}
      \right)^2,
    \label{eq:chi2}
\end{equation}
minimised analytically by:
\begin{equation}
    \theta_{12}^{\text{best}}
    = \arcsin\!\sqrt{P_{ee}^{\text{obs}}},
    \qquad
    \sigma_{\theta_{12}}
    = \frac{\sigma_{P_{ee}}}{\sin 2\theta_{12}}.
    \label{eq:bestfit}
\end{equation}
 
% ══════════════════════════════════════════════════════════════════════════════
\section{Results}
\label{sec:results}
 
\subsection{Best-Fit Mixing Angle}
 
Our central result is:
\begin{equation}
    \boxed{\theta_{12} = 36.0^{\circ} \pm 2.3^{\circ}
    \quad (\text{SNO Phase~I only}).}
    \label{eq:result}
\end{equation}
Table~\ref{tab:results} summarises all key quantities.
 
\begin{table}[h!]
    \centering
    \caption{Summary of experimental inputs and fit results.}
    \label{tab:results}
    \begin{tabular}{lcc}
        \toprule
        Quantity & Value & Units \\
        \midrule
        $\phi_\text{CC}$ & $1.76 \pm 0.11$
            & $10^6\ \text{cm}^{-2}\,\text{s}^{-1}$ \\
        $\phi_\text{NC}$ & $5.09 \pm 0.45$
            & $10^6\ \text{cm}^{-2}\,\text{s}^{-1}$ \\
        \midrule
        $P_{ee}^{\text{obs}}$ (CC/NC) & $0.3458 \pm 0.0374$ & --- \\
        \midrule
        $\theta_{12}$ (this work)     & $36.0^\circ \pm 2.3^\circ$ & --- \\
        $P_{ee}^{\text{MSW}}$ (fit)   & $0.346$ & --- \\
        $P_{ee}^{\text{vac}}$ at fit $\theta$ & $0.548$ & ruled out \\
        \midrule
        $\theta_{12}$ (PDG combined)  & $33.4^\circ \pm 0.8^\circ$ & --- \\
        \bottomrule
    \end{tabular}
\end{table}
 
\subsection{Panel A: Survival Probability Models}
 
\begin{figure}[h!]
    \centering
    \includegraphics[width=0.80\textwidth]{panel_A.png}
    \caption{%
        Electron-neutrino survival probability $P_{ee}$ as a function of
        $\theta_{12}$.
        Blue: adiabatic MSW prediction $\sin^2\theta_{12}$.
        Orange dashed: vacuum-averaged prediction
        $1 - \tfrac{1}{2}\sin^2 2\theta_{12}$.
        Green band: SNO-measured $P_{ee}^{\text{obs}} = 0.346 \pm 0.037$.
        Crimson dashed: PDG combined value $\theta_{12} = 33.4^\circ$.
    }
    \label{fig:panelA}
\end{figure}
 
Figure~\ref{fig:panelA} shows the most important conceptual result of
this analysis.
The green horizontal band is what SNO actually measured:
$P_{ee}^{\text{obs}} = 0.346 \pm 0.037$.
The blue MSW curve $\sin^2\theta_{12}$ passes cleanly through this band
at $\theta_{12} \approx 36^\circ$, consistent with our fit.
 
The orange vacuum curve $1 - \tfrac{1}{2}\sin^2 2\theta_{12}$ tells a
completely different story.
Its minimum value is $0.5$, reached at maximal mixing ($\theta_{12} = 45^\circ$).
It never reaches the observed level of $\sim\!35\%$ for any physical value
of $\theta_{12}$.
The vacuum formula is therefore not merely a poor fit --- it is
\emph{structurally incompatible} with the SNO measurement, regardless of
the angle chosen.
This is the graphical proof that matter effects dominate solar neutrino
propagation.
 
\subsection{Panel B: Chi-Squared Profile}
 
\begin{figure}[h!]
    \centering
    \includegraphics[width=0.80\textwidth]{panel_B.png}
    \caption{%
        Chi-squared profile $\Delta\chi^2(\theta_{12})$.
        Blue solid: MSW model.
        Purple dashed: vacuum model, shown for comparison.
        Shaded regions: $1\sigma$ ($\Delta\chi^2 < 1$) and
        $2\sigma$ ($\Delta\chi^2 < 4$) confidence intervals.
        Dark blue dotted: best-fit $36.0^\circ \pm 2.3^\circ$.
        Crimson dashed: PDG combined value $33.4^\circ$.
    }
    \label{fig:panelB}
\end{figure}
 
Figure~\ref{fig:panelB} quantifies the fit quality and uncertainty.
The MSW profile (blue) has a clean minimum at $\theta_{12} = 36.0^\circ$.
The $1\sigma$ confidence interval spans approximately
$[33.8^\circ,\ 38.3^\circ]$.
 
Two features deserve attention.
First, the crimson PDG line at $33.4^\circ$ falls just outside the
$1\sigma$ boundary at $\Delta\chi^2 \approx 1.3$.
This is not a contradiction --- it reflects the known systematic bias
in the SNO-only result that is corrected once the Earth's regeneration
effect and the KamLAND constraint are included, as discussed in
Section~\ref{sec:discussion}.
 
Second, the vacuum model's profile (purple dashed) never reaches zero ---
its minimum occurs near $45^\circ$, where $P_{ee}^{\text{vac}} = 0.5$,
still far above the observed value.
The two profiles are visually separated across the entire physical range,
confirming that the choice of model is not a matter of precision but of
correctness.
 
\subsection{Panel C: B8 Spectrum and Missing Neutrinos}
 
\begin{figure}[h!]
    \centering
    \includegraphics[width=0.80\textwidth]{panel_C.png}
    \caption{%
        Theoretical ${}^8\text{B}$ spectrum from the SSM (blue dashed)
        compared with the SNO-derived observed $\nu_e$ spectrum (black
        points).
        Orange shading: neutrinos that oscillated to $\nu_\mu/\nu_\tau$
        and were invisible to the CC channel.
        Green shading: surviving $\nu_e$ detected via CC.
        Red curve: best-fit MSW model at $\theta_{12} = 36.0^\circ$.
    }
    \label{fig:panelC}
\end{figure}
 
Figure~\ref{fig:panelC} presents the spectral picture of the solar
neutrino deficit.
The blue dashed curve is the theoretical ${}^8\text{B}$ flux shape
predicted by the SSM, normalised to unit area over the SNO-sensitive
range of $3.5$--$12$~MeV.
The black points show the observed CC-derived spectrum, obtained by
scaling the theoretical shape by $P_{ee}^{\text{obs}} = 0.346$.
 
The shaded regions make the physics immediately visible.
The \textbf{orange region} --- the gap between the SSM prediction and
the observed spectrum --- represents neutrinos that have oscillated to
$\nu_\mu$ or $\nu_\tau$.
These neutrinos are still present (detected via NC) but are invisible
to the CC channel.
The \textbf{green region} represents the surviving electron neutrinos
detected via CC interactions.
 
Crucially, the orange region is \emph{larger} than the green region,
meaning more than half the solar neutrinos have changed flavour by the
time they reach the detector.
This is the visual statement of MSW-dominated flavour conversion: the
Sun is producing predominantly $\ket{\nu_2}$ mass eigenstates, of which
only $\sin^2\theta_{12} \approx 35\%$ overlap with the $\nu_e$ flavour
state.
 
\subsection{Panel D: Numerical Summary}
 
\begin{figure}[h!]
    \centering
    \includegraphics[width=0.60\textwidth]{panel_D.png}
    \caption{%
        Numerical summary of the SNO Phase~I analysis: experimental
        inputs, derived survival probability, fit result, and comparison
        with the vacuum prediction and the PDG combined value.
    }
    \label{fig:panelD}
\end{figure}
 
Figure~\ref{fig:panelD} collects the numerical results for reference.
The measured fluxes $\phi_\text{CC}$ and $\phi_\text{NC}$ from
Reference~\cite{SNO2002} are the sole experimental inputs.
All other quantities are derived from these two numbers through the MSW
framework of Section~\ref{sec:theory}.
 
The ruling-out of the vacuum prediction is made quantitative here:
$P_{ee}^{\text{vac}} = 0.548$ at the best-fit angle, compared to the
observed $0.346$ --- a discrepancy of
$(0.548 - 0.346)/0.037 \approx 5.5\sigma$.
No value of $\theta_{12}$ brings the vacuum formula into agreement
with the data.
 
% ══════════════════════════════════════════════════════════════════════════════
\section{Discussion}
\label{sec:discussion}
 
\subsection{Why the Vacuum Formula Fails}
 
The vacuum-averaged survival probability has a hard lower bound of $0.5$
at maximal mixing.
SNO measured $P_{ee} \approx 0.346$, which lies below this bound.
No adjustment of $\theta_{12}$ can fix this --- the formula is
structurally wrong for this physical regime.
 
The reason is that the vacuum formula assumes the neutrino accumulates
a phase that averages to a $50/50$ mixture.
But ${}^8\text{B}$ neutrinos are produced deep inside the Sun where
matter effects rotate the mixing angle to $\theta_m \approx 90^\circ$.
The neutrino exits as a pure mass eigenstate $\ket{\nu_2}$, and the
overlap $|\braket{\nu_e}{\nu_2}|^2 = \sin^2\theta_{12}$ is set entirely
by the vacuum mixing angle --- not by any oscillation phase.
 
\subsection{Limitations of the SNO-Only Result}
 
Our result of $\theta_{12} = 36.0^\circ \pm 2.3^\circ$ has two notable
limitations relative to the PDG value of $33.4^\circ \pm 0.8^\circ$.
 
\paragraph{Statistical uncertainty ($\pm 2.3^\circ$).}
The SNO CC/NC measurement provides only a single integrated ratio with
no energy resolution.
The adiabatic MSW formula $P_{ee} = \sin^2\theta_{12}$ is independent
of $\Delta m^2_{21}$, so SNO alone cannot constrain the mass-squared
splitting.
This leaves a parameter degeneracy that inflates the uncertainty by
a factor of $\sim\!3$ compared to the combined fit.
 
\paragraph{Systematic offset ($\sim\!2.6^\circ$).}
Two effects bias the SNO-only result upward.
 
\subsection{The Earth's Day-Night Regeneration Effect}
 
When solar neutrinos pass through the Earth at night, they encounter
increasing electron density.
This triggers a partial reverse MSW conversion: some $\nu_\mu/\nu_\tau$
flux partially reconverts to $\nu_e$, slightly elevating the nighttime
CC rate.
Without correcting for this, the time-averaged $P_{ee}^{\text{obs}}$
is slightly overestimated.
The corrected survival probability is:
\begin{equation}
    P_{ee}^{\text{corrected}}
    = \sin^2\theta_{12} + \delta(N_e,\, \Delta m^2_{21},\, E),
    \label{eq:regen}
\end{equation}
where $\delta$ is computed by integrating the MSW Hamiltonian through
the Earth's density profile (PREM \cite{PREM}).
Including this term pulls the fitted $\theta_{12}$ downward toward
the true value.
 
\subsection{The KamLAND Constraint}
 
KamLAND \cite{KamLAND2003} measured $\bar\nu_e$ disappearance from
Japanese reactors at $L \sim 180$~km, observing an energy-resolved
oscillation pattern:
\begin{equation}
    P_{\bar\nu_e \to \bar\nu_e}(L, E)
    = 1 - \sin^2 2\theta_{12}\,
      \sin^2\!\left(\frac{1.27\,\Delta m^2_{21}\,L}{E}\right).
    \label{eq:kamland}
\end{equation}
This provided a precise independent measurement of $\Delta m^2_{21}$.
Combined with SNO in a global fit \cite{PDG2022}:
\begin{equation}
    \theta_{12} = 33.4^\circ \pm 0.8^\circ,
    \qquad
    \Delta m^2_{21} = (7.53 \pm 0.18)\times 10^{-5}\ \text{eV}^2.
    \label{eq:pdg}
\end{equation}
The combination broke the parameter degeneracy, reduced the uncertainty
by a factor of $\sim\!3$, and pulled the central value from $36.0^\circ$
to $33.4^\circ$.
 
% ══════════════════════════════════════════════════════════════════════════════
\section{Conclusion}
\label{sec:conclusion}
 
We have reproduced the SNO Phase~I measurement of the solar neutrino
mixing angle using the CC/NC flux ratio as the primary observable.
Working in the physically correct adiabatic MSW regime, we obtain
\[
    \theta_{12} = 36.0^{\circ} \pm 2.3^{\circ},
\]
consistent with the original SNO-only result published in 2002.
 
The analysis demonstrates three key points:
\begin{enumerate}
    \item The vacuum-averaged oscillation formula is \emph{structurally
          incompatible} with SNO data --- no value of $\theta_{12}$
          brings $P_{ee}^{\text{vac}} \geq 0.5$ into agreement with
          the observed $P_{ee} \approx 0.35$.
          Matter effects are not a perturbative correction; they are the
          dominant physics.
 
    \item The adiabatic MSW limit yields $P_{ee} = \sin^2\theta_{12}$,
          a direct consequence of the neutrino exiting the Sun as the
          pure mass eigenstate $\ket{\nu_2}$.
 
    \item The SNO-only offset from the PDG value is physically understood:
          it arises from neglecting the Earth's day-night regeneration
          effect and the $\Delta m^2_{21}$ constraint subsequently
          provided by KamLAND.
\end{enumerate}
 
The story of solar neutrino oscillations stands as one of the most
compelling in modern physics: a deficit first observed by Davis in 1968
\cite{Homestake1968}, explained by Mikheyev, Smirnov, and Wolfenstein
in the 1980s \cite{MSW1985, MSW1986}, and confirmed by SNO and KamLAND
in the early 2000s --- earning the 2015 Nobel Prize in Physics shared
by Takaaki Kajita and Arthur B.~McDonald.
 
% ══════════════════════════════════════════════════════════════════════════════
\begin{thebibliography}{99}
 
\bibitem{Homestake1968}
    R.~Davis~Jr., D.~S.~Harmer, and K.~C.~Hoffman,
    ``Search for neutrinos from the sun,''
    \textit{Phys.\ Rev.\ Lett.}\ \textbf{20}, 1205 (1968).
 
\bibitem{Bahcall1989}
    J.~N.~Bahcall,
    \textit{Neutrino Astrophysics}.
    Cambridge University Press, 1989.
 
\bibitem{MSW1985}
    S.~P.~Mikheyev and A.~Yu.~Smirnov,
    ``Resonance enhancement of oscillations in matter and solar neutrino
    spectroscopy,''
    \textit{Sov.\ J.\ Nucl.\ Phys.}\ \textbf{42}, 913 (1985).
 
\bibitem{MSW1986}
    L.~Wolfenstein,
    ``Neutrino oscillations in matter,''
    \textit{Phys.\ Rev.\ D}\ \textbf{17}, 2369 (1978);
    \textit{ibid.}\ \textbf{20}, 2634 (1979).
 
\bibitem{SNO2002}
    Q.~R.~Ahmad \textit{et al.}\ (SNO Collaboration),
    ``Measurement of day and night neutrino energy spectra at the
    Sudbury Neutrino Observatory,''
    \textit{Phys.\ Rev.\ Lett.}\ \textbf{89}, 011302 (2002).
    \href{https://arxiv.org/abs/nucl-ex/0204009}{arXiv:nucl-ex/0204009}.
 
\bibitem{KamLAND2003}
    K.~Eguchi \textit{et al.}\ (KamLAND Collaboration),
    ``First results from KamLAND: Evidence for reactor antineutrino
    disappearance,''
    \textit{Phys.\ Rev.\ Lett.}\ \textbf{90}, 021802 (2003).
    \href{https://arxiv.org/abs/hep-ex/0212021}{arXiv:hep-ex/0212021}.
 
\bibitem{PREM}
    A.~M.~Dziewonski and D.~L.~Anderson,
    ``Preliminary reference Earth model,''
    \textit{Phys.\ Earth Planet.\ Inter.}\ \textbf{25}, 297 (1981).
 
\bibitem{PDG2022}
    R.~L.~Workman \textit{et al.}\ (Particle Data Group),
    ``Review of Particle Physics,''
    \textit{Prog.\ Theor.\ Exp.\ Phys.}\ \textbf{2022}, 083C01 (2022).
 
\end{thebibliography}
 
\end{document}